% !TeX root = ../../main.tex
\subsection{Training Stages and Checkpoint Selection}
The continuation experiments start from the validation-selected partial-map
P/G checkpoint, inherited from an earlier outcome-supervised mobile policy.
This two-level H-JEPA model already contains rank-8 LoRA and P/G outcome heads
adapted to simulated door opening. We resume the adapters with scaling $\alpha=16$ and
dropout 0.05 \citep{hu2021lora}. Table~\ref{tab:training} gives the latest
continuation budget. AdamW uses weight decay 0.01 and gradient clipping at 1;
P/G training uses BF16 autocast. Groups contain four related layouts, and the
paired sampler draws two different layouts from one training group per update.

\begin{table}[H]
\caption{Latest continuation settings. Learning rates are for the standard arm;
the low-rate controls multiply both P/G-stage rates by 0.1.}
\label{tab:training}
\centering\small
\begin{tabularx}{\linewidth}{@{}lXXX@{}}
\toprule
Stage & Trainable parameters & Updates and batch & Learning rate \\
\midrule
P/G & Existing LoRA and P/G heads & 320; 2 scenes; target refresh every 40
 & Heads \mbox{$3\times10^{-4}$}\newline LoRA \mbox{$2\times10^{-5}$} \\
I & Inspection head & 1600; up to 32 training rows & $3\times10^{-4}$ \\
\bottomrule
\end{tabularx}
\end{table}

P/G checkpoint selection compares direct and always-inspect validation
branches: fewer collisions first, then inspected-branch correctness, direct
correctness, total openings, net utility, and path length. Strict improvement
retains the earlier checkpoint on ties. We then freeze all non-I parameters
and select I using validation task outcomes and cost. Python, NumPy, Torch,
and CUDA use seed 42.
